% CV de Yacouba BARRO — version française
% Compiler avec : xelatex cv-fr.tex
\documentclass[10pt]{article}
\usepackage{cv-style}

\begin{document}

\cvheader{Yacouba}{BARRO}
  {Élève ingénieur statisticien-économiste (ISE)}
  {Ouagadougou, Burkina Faso\sep (+226) 66~54~84~30 / 68~54~36~46\\[1pt]\href{mailto:barroyacouba192962@gmail.com}{barroyacouba192962@gmail.com}}
  {\href{https://bdigistat.com}{bdigistat.com}\sep \href{https://github.com/barroyacouba192962-sys}{github.com/barroyacouba192962-sys}\sep \href{https://www.linkedin.com/in/barro-undefined-b30b47340/}{LinkedIn}}
  {photo.png}

\section{Profil}
Élève ingénieur statisticien-économiste à l'ISSP (Université Joseph Ki-Zerbo) et titulaire d'une licence en économie,
je me spécialise en économétrie appliquée, en macroéconomie quantitative et en science des données. Je m'intéresse à la modélisation macroéconomique, à l'économétrie appliquée aux économies ouest-africaines et à
l'apport de l'apprentissage automatique à la prévision. En parallèle, je conçois des outils numériques et des formations
qui rendent les méthodes quantitatives accessibles au Burkina Faso. Je recherche un stage, une bourse ou un projet
doctoral dans ces domaines.

\section{Centres d’intérêt de recherche}
Macroéconomie appliquée et modèles DSGE\sep Économétrie des séries temporelles\sep Commerce, investissements
directs étrangers et croissance en Afrique de l'Ouest\sep Prévision et apprentissage automatique\sep Économie du
développement et approche par les capabilités

\section{Formation}
\cventry{2023 – auj.}{Cycle d'Ingénieur Statisticien Économiste (ISE)}
  {Institut Supérieur des Sciences de la Population (ISSP), Université Joseph Ki-Zerbo, Ouagadougou}
  {Enseignements~: économétrie, comptabilité nationale (SCN~2008), macroéconomie, apprentissage automatique appliqué.
   Stage de fin de cycle en cours chez Coris Assurance Vie.}
\cventry{2023}{Licence en économie, option Macroéconomie et gestion du développement}
  {Université Nazi Boni, Bobo-Dioulasso}{}
\cventry{2018}{Baccalauréat}{Complexe Scolaire Alpha la Solidarité}{}

\section{Travaux de recherche et projets académiques}
\cventry{ISSP}{Modèle DSGE de type RBC appliqué au Burkina Faso (données 2005--2021)}{}
  {Construction et simulation d'un modèle d'équilibre général dynamique stochastique sur données burkinabè avec
   Dynare et Octave~; restitution sous forme de présentation Beamer (\LaTeX).}
\cventry{ISSP}{Ouverture commerciale, IDE et croissance économique au Burkina Faso}{}
  {Estimation d'un modèle ARDL (relations de court et de long terme) sous Python~; rédaction du rapport économétrique.}
\cventry{ISSP}{Prévision de la demande horaire d'électricité (série AEP, 2004--2018)}{}
  {Comparaison de LightGBM, XGBoost, forêts aléatoires, LSTM (PyTorch) et régression linéaire.
   Meilleur modèle~: LightGBM (MAPE~=~0,70\,\%, R²~=~0,997).}

\section{Expérience}
\cventry{En cours}{Stagiaire de fin de cycle}{Coris Assurance Vie, Burkina Faso}{}
\cventry{En cours}{Fondateur et développeur, Faso Carrières}
  {\href{https://fasocarrieres.com}{fasocarrieres.com} — plateforme d'études, de formation et d'emploi au Burkina Faso}
  {\begin{cvlist}
     \item Conception et développement de la plateforme (Next.js, React, Node.js, Appwrite, Vercel).
     \item Rédaction de formations en ligne~: collecte de données mobile (KoboToolbox, ODK), enquêtes et échantillonnage,
           économétrie appliquée (Stata, R), analyse de données (SPSS, Stata), cartographie (QGIS), \LaTeX{},
           Excel appliqué à la finance, IA générative.
     \item Projet en compétition au Concours National d'Innovation Digitale 2026 (Le Faso Digital).
   \end{cvlist}}
\cventry{1 mois}{Stagiaire}{Institut National de la Statistique et de la Démographie (INSD), Ouagadougou}{}
\cventry{UJKZ}{Assistant numérique}{Université Joseph Ki-Zerbo, Ouagadougou}
  {Mise à jour des données de la plateforme \href{https://campusfaso.com}{campusfaso.com}.}
\cventry{Indépendant}{Graphiste et vidéaste}{}
  {Étiquettes et packaging pour des marques de cosmétiques, affiches, supports de communication et montage vidéo.}

\section{Réalisations numériques et pédagogiques}
\cventry{Mobile}{FasoPrepas, application de préparation aux concours}{}
  {QCM chronométrés, corrections détaillées et suivi de progression pour les concours ISE, ASECNA, IFORD, etc.
   (Flutter, Node.js, Appwrite).}
\cventry{Édition}{La RévueCapes}{}
  {Mise en forme et édition d'un document de préparation au CAPES option mathématiques (\LaTeX).}

\section{Certifications}
\cventry{Coursera}{Finance \& Quantitative Modeling for Analysts (spécialisation)}{University of Pennsylvania}
  {\href{https://coursera.org/verify/specialization/R352N0T1ED2I}{coursera.org/verify/specialization/R352N0T1ED2I}}
\cventry{Coursera}{Microsoft Power BI Data Analyst}{Microsoft}
  {\href{https://coursera.org/verify/23P9VXRLDMQJ}{coursera.org/verify/23P9VXRLDMQJ}}
\cventry{Autre}{Formation en infographie}{}{}

\section{Compétences}
\cvitem{Économétrie}{Stata, R, SPSS, Python (scikit-learn, LightGBM, XGBoost, PyTorch), Dynare, Octave}
\cvitem{Données}{KoboToolbox, ODK, QGIS, Power BI (Power Query), Excel (VBA), SQL, PostgreSQL, Access}
\cvitem{Rédaction}{\LaTeX{} et Beamer, Word}
\cvitem{Web \& mobile}{Next.js, React, Node.js, Flutter (Dart), Appwrite, Git et GitHub}
\cvitem{Design}{Photoshop, Illustrator, InDesign, Canva, CapCut}

\section{Langues}
\cvitem{Maternelles}{Dioula, Bolon}
\cvitem{Français}{Courant (langue de travail)}
\cvitem{Anglais}{Niveau de base, lecture de documentation technique}

\section{Engagement et centres d’intérêt}
\cvitem{Associatif}{N'Dorola Massakô~: conception d'affiches de collecte de fonds et de supports de communication}
\cvitem{Intérêts}{Économie politique et philosophie du développement (Amartya Sen), éducation, innovation numérique,
  création de contenus éducatifs}
\cvitem{Références}{Disponibles sur demande}

\end{document}
